
\documentclass[a4paper,10pt]{report}

%\usepackage[landscape]{geometry}
\usepackage[cm-default]{fontspec}
\setromanfont{FreeSerif}
\setsansfont{FreeSans}
\setmonofont{FreeMono}
\usepackage[utf8x]{inputenc}
\usepackage{fontspec}
\usepackage{xunicode}
\usepackage{xltxtra}
\usepackage{xgreek}
\usepackage{amsmath}
\usepackage{unicode-math}
\usepackage{ulem}
\usepackage{color}
\usepackage{verbatim}
\usepackage{nopageno}
\usepackage{graphicx}
\usepackage{textpos}
\setlength{\TPHorizModule}{1cm}
\setlength{\TPVertModule}{1cm}
%\usepackage[colorgrid,texcoord]{eso-pic}
\usepackage[outline]{contour}
\usepackage{wrapfig}
\usepackage{url}
\usepackage{color}
\usepackage{tikz}
\usepackage{fancybox,fancyhdr}
\usepackage{subfigure}
\usepackage{pstricks}
\usepackage{epsfig}
\usepackage{multicol}
\usepackage{listings}
\usepackage{enumerate}
\usepackage{hyperref}
\hypersetup{
  bookmarks=true,
  bookmarksopen=true,
  pdfborder=false,
  pdfpagemode=UseNone,
  raiselinks=true,
  pdfhighlight={/P},
  colorlinks,
  citecolor=black,
  filecolor=black,
  linkcolor=black,
  urlcolor=black
}


\usepackage{fontspec}
\usepackage{xunicode}
\usepackage{xltxtra}


% Margins
%\setlength{\textwidth}{24cm}

%\setlength{\voffset}{0in}
%\setlength{\textheight}{6.8in}
% Colors
\definecolor{cbrown}{rgb}{0.49,0.24,0.07}
\definecolor{cmpez}{rgb}{0.92,0.88,0.79}
\definecolor{cmpez2}{rgb}{0.62,0.33,0.34}
\definecolor{cwhite}{rgb}{1,1,1}
\definecolor{cblack}{rgb}{0,0,0}
\definecolor{cred}{rgb}{0.9,0.15,0.15}
\definecolor{clblue}{rgb}{0.57,0.85,0.97}
\definecolor{corange}{rgb}{0.86,0.49,0.18}
\definecolor{clorange}{rgb}{0.95,0.84,0.61}
\definecolor{cyellow}{rgb}{1,0.95,0.16}
\definecolor{cgreen}{rgb}{0.69,0.8,0.31}
\definecolor{clbrown}{rgb}{0.78,0.67,0.43}
\definecolor{cmagenta}{rgb}{0.79,0.05,0.54}
\definecolor{cgray}{rgb}{0.76,0.73,0.63}
\definecolor{bluesite}{rgb}{0.05,0.43,0.67}
\definecolor{redsite}{rgb}{0.90,0.15,0.15}
\definecolor{backsite}{rgb}{0.07,0.07,0.07}
% ----------- from Costas -----------------------------

%---------------   Main Fonts    -----------------%   

% \setmainfont[Mapping=tex-text]{Calibri}
% \setmainfont[Mapping=tex-text]{Times New Roman}
%\setmainfont[Mapping=tex-text]{Droid Serif}
%\setmainfont[Mapping=tex-text]{Cambria}
%\setmainfont[Mapping=tex-text]{cm-unicode}
% \setmainfont[Mapping=tex-text]{Gentium}
% \setmainfont[Mapping=tex-text]{GFS Didot}
% \setmainfont[Mapping=tex-text]{Comic Sans MS}
 \setmainfont[Mapping=tex-text]{Ubuntu}
% \setmainfont[Mapping=tex-text]{Myriad Pro}
% \setmainfont[Mapping=tex-text]{CMU Concrete}
%\setmainfont[Mapping=tex-text]{DejaVu Sans}
%\setmainfont[Mapping=tex-text]{KerkisSans}
%\setmainfont[Mapping=tex-text]{KerkisCaligraphic}
% \setmainfont[Mapping=tex-text]{Segoe Print}
% \setmainfont[Mapping=tex-text]{Gabriola}

%------------------------------------------------%

% ------------ Mathematics Fonts ----------------%
\setmathfont{Asana-Math.ttf}
% -----------------------------------------------%
% Misc
\author{Κασωτάκη Ε. - Σμαραγδάκης Κ.}
\title{www.math24.gr}
% ---------------- Formation --------------------%


\setlength\topmargin{-0.7cm}
\addtolength\topmargin{-\headheight}
\addtolength\topmargin{-\headsep}
\setlength\textheight{26cm}
\setlength\oddsidemargin{-0.54cm}
\setlength\evensidemargin{-0.54cm}
%\setlength\marginparwidth{1.5in}
\setlength\textwidth{17cm}

\RequirePackage[avantgarde]{quotchap}
\renewcommand\chapterheadstartvskip{\vspace*{0\baselineskip}}
\RequirePackage[calcwidth]{titlesec}
\titleformat{\section}[hang]{\bfseries}
{\Large\thesection}{12pt}{\Large}[{\titlerule[0.9pt]}]
%--------------------------------------------------%
\usepackage{draftwatermark}
\SetWatermarkText{www.math24.gr}
\SetWatermarkLightness{0.9}
\SetWatermarkScale{0.6}
%--------------------------------------------------%


\usepackage{xcolor}
\usepackage{amsthm}
\usepackage{framed}
\usepackage{parskip}

\colorlet{shadecolor}{bluesite!20}

\newtheorem*{orismos}{Ορισμός}

\newenvironment{notation}
  {\begin{shaded}\begin{theorem}}
  {\end{theorem}\end{shaded}}




% Document begins
\begin{document}
\pagestyle{fancy}
\fancyhead{}
\fancyfoot{}
\renewcommand{\headrulewidth}{0pt}
\renewcommand{\footrulewidth}{0pt}


\fancyhead[LO,LE]{
 %\textblockcolor{backsite}
 %\begin{textblock}{5}(-2,-0.55)
  %\rule{0cm}{1cm}
 %\end{textblock}
 \textblockcolor{bluesite}
 \begin{textblock}{5}(-1.5,-0.55)
  \rule{0cm}{1cm}
 \end{textblock}
 %\textblockcolor{bluesite}
 %\begin{textblock}{14}(5.5,-0.55)
 % \rule{0cm}{1cm}
 %\end{textblock}
 \begin{textblock}{0}(-1,-0.25)
 \color{cwhite} \begin{Large}www.math24.gr\end{Large}
 \end{textblock}
\begin{textblock}{0}(16,-0.4)
 \color{cwhite} \includegraphics[height=1.5cm]{math24_logo.png}
 \end{textblock}
\textblockcolor{white}
\begin{textblock}{17}(-1,28)
 \color{backsite} \begin{small}Copyright \textcopyright 2011, Κασωτάκη Ε.(ikasotaki@gmail.com) - Σμαραγδάκης Κ.(kesmarag@gmail.com)\end{small}
 \end{textblock}
}
 \begin{shaded}
\begin{center}
\huge \textbf{Μαθηματικά Α' Γυμνασίου}\\
\end{center} 
\begin{center}
\textbf{"Επαναληπτικές Ασκήσεις για τις διακοπές του Πάσχα στην Άλγεβρα"} \hfill \textbf{}
\end{center}
%\subsection*{}
\end{shaded}
%\vspace{2em
%\begin{shaded}
%\begin{center}
%\huge \textbf{Επαναληπτικές Ασκήσεις για το Κεφάλαιο 1}\\
%Πρόσθεση ρητών αριθμών
%\end{center} 
%\textbf{Μαθηματικά Α' Γυμνασίου} \hfill \textbf{Ημερομηνία Παράδοσης : \hspace{2em} }
%\subsection*{Ονοματεπώνυμο :\hfill  \hspace{5em}}
%\end{shaded}
\vspace{2em}
\begin{itemize}
\item Κεφάλαιο 1: Φυσικοί αριθμοί 
\item Κεφάλαιο 2: Κλάσματα
\item Κεφάλαιο 3: Δεκαδικοί αριθμοί 
\item Κεφάλαιο 4: Εξισώσεις
\item Κεφάλαιο 7: Θετικοί και αρνητικοί αριθμοί
\end{itemize}


\section*{Άσκηση 1  \hfill \small{}}
Να βρείτε το $ΕΚΠ$ και το $ΜΚΔ$  των παρακάτω αριθμών:
\begin{multicols}{2}
\begin{enumerate}[1)]
 \item 15 και 20
 \item 11 και 22
 \item 12 και 18
 \item 14 και 34
 \item 21 και 44
 \item 111 και 444
 \item 183 και 346
 \item 225 και 730
 \item 500 και 5000
 \item 128 και 1024
\end{enumerate}
\end{multicols}


\section*{Άσκηση 2  \hfill \small{}}
Για κάθε ένα από τους παρακάτω αριθμούς να γράψετε το ανάπτυγμα  με χρήση των δυνάμεων του 10.
\begin{multicols}{2}
\begin{enumerate}[1)]
 \item 8231
 \item 1345
 \item 2628
 \item 3344
 \item 6782
 \item 13435
 \item 37373
 \item 44444
 \item 50000
 \item 65871
\end{enumerate}
\end{multicols}



\section*{Άσκηση 3  \hfill \small{}}
Να γράψετε ένα \textbf{ναι} ή ένα \textbf{όχι} στις κενές θέσεις του παρακάτω πίνακα σύμφωνα με το εάν οι αριθμοί 
που δίνονται διαιρούνται με τους 2,3,4,5,9,19,25 και 100.
\begin{center}
\begin{tabular}{|c|c|c|c|c|c|c|c|c|}
\hline
             & \textbf{2} & \textbf{3} & \textbf{4} & \textbf{5} & \textbf{9} & \textbf{10} & \textbf{25} & \textbf{100}\\
\hline
\textbf{100} &  &  &  &  &  &  &  & \\
\hline
\textbf{111} &  &  &  &  &  &  &  & \\
\hline
\textbf{112} &  &  &  &  &  &  &  & \\
\hline
\textbf{313} &  &  &  &  &  &  &  & \\
\hline
\textbf{414} &  &  &  &  &  &  &  & \\
\hline
\textbf{650} &  &  &  &  &  &  &  & \\
\hline
\textbf{7213} & &  &  &  & &  &  & \\
\hline
\textbf{8291} & &  & &  &  &  &  & \\
\hline
\textbf{9225} & &  &  &  &  &  &  & \\
\hline
\textbf{9999} & &  &  & &  &  &  & \\
\hline
\end{tabular}
\end{center}



\section*{Άσκηση 4  \hfill \small{}}
Να απλοποιήσετε τα παρακάτω κλάσματα:
\begin{multicols}{2}
\begin{enumerate}[1)]
 \item $\dfrac{4}{16}$
 \item $\dfrac{81}{42}$
 \item $\dfrac{120}{66}$
 \item $\dfrac{4\cdot 3\cdot 5}{15\cdot16\cdot12}$
 \item $\dfrac{3\cdot 22\cdot 100}{4\cdot 44\cdot 200}$
 \item $\dfrac{2^{3}\cdot4^{3}\cdot5^{2}\cdot25}{8\cdot100\cdot16}$
 \item $\dfrac{2^{3}\cdot3^{4}\cdot5^{6}}{2^{2}\cdot3^{3}\cdot5^{5}}$
 \item $\dfrac{12\cdot18\cdot20}{6\cdot4\cdot5}$
 \item $\dfrac{7^{2}\cdot21\cdot3^{4}}{7^{3}\cdot3^{5}}$
 \item $\dfrac{25\cdot30\cdot2^{2}}{5^{2}\cdot20}$
\end{enumerate}
\end{multicols}



\section*{Άσκηση 5  \hfill \small{}}
Να εξετάσετε εάν τα παρακάτω κλάσματα είναι ισοδύναμα:
\begin{multicols}{2}
\begin{enumerate}[1)]
 \item $\dfrac{1}{2}$ και $\dfrac{1}{3}$
 \item $\dfrac{3}{4}$ και $\dfrac{6}{8}$
 \item $\dfrac{7}{2}$ και $\dfrac{2}{7}$
 \item $\dfrac{8}{3}$ και $\dfrac{4}{6}$
 \item $\dfrac{2}{11}$ και $\dfrac{4}{22}$
 \item $\dfrac{10}{30}$ και $\dfrac{10}{40}$
 \item $\dfrac{10}{30}$ και $\dfrac{2}{15}$
 \item $\dfrac{3}{16}$ και $\dfrac{6}{32}$
 \item $\dfrac{3}{12}$ και $\dfrac{1}{4}$
 \item $\dfrac{8}{80}$ και $\dfrac{1}{10}$
\end{enumerate}
\end{multicols}



\section*{Άσκηση 6  \hfill \small{}}
Να βρείτε τους αντίστροφους των παρακάτω αριθμών:
\begin{multicols}{2}
\begin{enumerate}[1)]
 \item $13$
 \item $\dfrac{2}{7}$
 \item $\dfrac{5}{3}$
 \item $\dfrac{1}{6}$
 \item $\dfrac{2}{20}$
 \item $117$
 \item $\dfrac{13}{27}$
 \item $\dfrac{81}{3}$
 \item $\dfrac{4}{32}$
 \item $1$
\end{enumerate}
\end{multicols}

\section*{Άσκηση 7  \hfill \small{}}
Να γράψετε τα παρακάτω κλάσματα ως δεκαδικούς αριθμούς:
\begin{multicols}{2}
\begin{enumerate}[1)]
 \item $\dfrac{71}{10}$
 \item $\dfrac{88}{100}$
 \item $\dfrac{3}{1000}$
 \item $\dfrac{27}{10}$
 \item $\dfrac{600}{10}$
 \item $\dfrac{20}{4}$
 \item $\dfrac{11}{2}$
 \item $\dfrac{4}{25}$
 \item $\dfrac{1}{5}$
 \item $\dfrac{36}{6}$
\end{enumerate}
\end{multicols}



\section*{Άσκηση 8  \hfill \small{}}
Να γράψετε ως δεκαδικά κλάσματα κάθε ένα από τους παρακάτω δεκαδικούς αριθμούς:
\begin{multicols}{2}
\begin{enumerate}[1)]
 \item $1,1$
 \item $0,91$
 \item $7,01$
 \item $0,133$
 \item $78,2$
 \item $41,41$
 \item $0,007$
 \item $8,088$
 \item $0,096$
 \item $71,006$
\end{enumerate}
\end{multicols}


\section*{Άσκηση 9  \hfill \small{}}
Για κάθε ένα από τους παρακάτω αριθμούς να βρείτε τον αντίθετο τους:
\begin{multicols}{2}
\begin{enumerate}[1)]
 \item $4$
 \item $\dfrac{2}{3}$
 \item $-\dfrac{4}{5}$
 \item $-\dfrac{1}{6}$
 \item $\dfrac{4}{9}$
 \item $\dfrac{8}{16}$
 \item $9,54$
 \item $-3$
 \item $-22,81$
 \item $81$
\end{enumerate}
\end{multicols}

\section*{Άσκηση 10  \hfill \small{}}
Να υπολογίσετε το αποτέλεσμα για κάθε μία από τις παρακάτω αριθμητικές παραστάσεις εφαρμόζοντας πρώτα την 
επιμεριστική ιδιότητα:
%\begin{multicols}{2}
\begin{enumerate}[1)]
 \item $2\cdot (3+7+\dfrac{1}{2}+\dfrac{6}{4}-\dfrac{3}{2})$
 \item $10\cdot (8,1-0,1-40-\dfrac{7}{10}+\dfrac{40}{100})$
 \item $100\cdot (-0,02+0,01+\dfrac{73}{100}-\dfrac{80}{100})$
 \item $4\cdot(0,25+\dfrac{8}{4}+\dfrac{1}{16})-\dfrac{8}{32}+
        2\cdot(-0,05+\dfrac{1}{2}+4^{2}-\dfrac{32}{2})$
 \item $3\cdot (\dfrac{11}{3}+|2|-\dfrac{5}{3})+\dfrac{2}{6}
       -|-2|\cdot(\dfrac{8}{3}-\dfrac{8}{12}+\dfrac{1}{2})$
 \item $2^{2}\cdot(3-7+\dfrac{8}{12})-3^{2}\cdot(0-1+\dfrac{1}{9})$
 \item $\dfrac{1}{2}\cdot(\dfrac{4}{2}-\dfrac{8}{6}-\dfrac{1}{4})
       +\dfrac{1}{3}\cdot(\dfrac{2}{4}-\dfrac{3}{2})$
 \item $6\cdot(1-\dfrac{36}{3}-\dfrac{18}{2})-|3|\cdot(|-1|+|0|+|-7|)$
 \item $2\cdot (\dfrac{4}{2}+\dfrac{3}{2}+\dfrac{2}{2})
       -3\cdot(\dfrac{1}{3}+\dfrac{2}{3}+\dfrac{4}{3})$
 \item $4\cdot (-\dfrac{6}{2^{2}}-\dfrac{4}{2^{2}})
        +3^{2}\cdot(\dfrac{9}{27}-\dfrac{6}{27})$
\end{enumerate}
%\end{multicols}

\section*{Άσκηση 11  \hfill \small{}}
Να βρείτε το αποτέλεσμα των παρακάτω αριθμητικών παραστάσεων χρησιμοποιώντας το αντίστροφο 
της επιμεριστικής ιδιότητας:
%\begin{multicols}{2}
\begin{enumerate}[1)]
 \item $7\cdot 3+3\cdot 2+3$
 \item $0,72\cdot10+0,72\cdot30-0,72\cdot4-0,72\cdot6-0,72\cdot70$
 \item $\dfrac{1}{2}\cdot3+2\cdot \dfrac{1}{2}+\dfrac{1}{2}-20\cdot \dfrac{1}{2}+10\cdot \dfrac{1}{2}$
 \item $3,68\cdot \dfrac{2}{10}+3,68\cdot \dfrac{3}{10}+3,68\cdot \dfrac{5}{10}
        -3,68\cdot \dfrac{3}{10}-3,68\cdot \dfrac{4}{10}+3,68\cdot \dfrac{7}{10}$
 \item $5\cdot |-4|+5\cdot |-3|+5\cdot |+2|-5\cdot |-3|+5\cdot |-1|$
 \item $\dfrac{1}{2}\cdot \dfrac{50}{4}-\dfrac{1}{2}\cdot \dfrac{60}{5}+\dfrac{1}{2}\cdot \dfrac{15}{10}$
\end{enumerate}
%\end{multicols}


\section*{Άσκηση 12  \hfill \small{}}
Να βάλετε σε αύξουσα σειρά τους παρακάτω αριθμούς:
%\begin{multicols}{2}
\begin{enumerate}[1)]
 \item $-7, 3,-3,0,9,20$
 \item $|-7|,-9,-4,|8|,21$
 \item $\dfrac{5}{2},-\dfrac{1}{2},\dfrac{4}{2},-\dfrac{7}{2},\dfrac{9}{2}$
 \item $-\dfrac{6}{25},\dfrac{6}{9},-\frac{6}{7},\dfrac{6}{5},\frac{6}{41}$
 \item $-\dfrac{6}{4},-\dfrac{2}{5},\dfrac{15}{10},\dfrac{70}{20},\dfrac{21}{2}$
 \item $-7,1\quad,\quad-7,2\quad,\quad 0 \quad ,\quad 7,6 \quad , \quad 7,7$
 \item $11,81\quad,\quad-13,71\quad,\quad 0 \quad ,\quad -47,2 \quad , \quad 7,8$
 \item $2^{2},-3^{2},-9^{2},3^{3},4^{2}$
 \item $8,6\quad , \quad 0, -\dfrac{1}{2},|-70|,-70$
 \item $-\dfrac{1}{3},-\dfrac{5}{4},\dfrac{3}{2},\dfrac{6}{6},0$
\end{enumerate}
%\end{multicols}


\section*{Άσκηση 13  \hfill \small{}}
Να τοποθετήσετε στην ευθεία των αριθμών τους παρακάτω αριθμούς:
%\begin{multicols}{2}
\begin{enumerate}[1)]
 \item $0,-5,4$
 \item $-\dfrac{1}{4},-1,0,\dfrac{1}{2}$
 \item $\dfrac{5}{2},-2,2,\dfrac{3}{2}$
 \item $\dfrac{3}{2},\dfrac{3}{2},\dfrac{5}{2},\dfrac{1}{2}$
 \item $-|2|,-|-3|,|7|,|-1|$
\end{enumerate}
%\end{multicols}

\section*{Άσκηση 14  \hfill \small{}}
Να υπολογίσετε τις τιμές των παρακάτω αριθμητικών παραστάσεων:
%\begin{multicols}{2}
\begin{enumerate}[1)]
 \item $2\cdot (4-3-7+6)+71\cdot(81-80-22+4+6+2)$
 \item $3\cdot \dfrac{7}{3}+3,1\cdot 10-2\cdot 2^{3}-1-81$
 \item $\dfrac{4}{3}\cdot \dfrac{5}{2}-\dfrac{1}{2}\cdot (3-2)+|-1|$
 \item $(7-3)\cdot(8,1-0,1)+(-7,8-1,1)\cdot\dfrac{1}{2}$
 \item $\dfrac{3}{2}\cdot \dfrac{2^{2}}{3^{2}}\cdot (-5)-(-\dfrac{1}{3})\cdot (\dfrac{6}{2^{2}})$
 \item $710\cdot(\dfrac{8}{10}-\dfrac{2}{5})+7,1\cdot(-|-7|+|-107|)-3\cdot(-78+8+40)$
 \item $\dfrac{6}{2}:\dfrac{4}{3}-\dfrac{1}{2}:(-\dfrac{2}{3})+(51-6+5)\cdot \dfrac{1}{2\cdot 25}$
 \item $|-1|\cdot|-8|\cdot |7|-\dfrac{1}{2}\cdot 100+1,8\cdot 2$
 \item $\dfrac{15}{12}\cdot \dfrac{3^{2}}{25}\cdot \dfrac{44}{4\cdot 5}$
 \item $(678\cdot 1)\cdot (-56+56)$
\end{enumerate}
%\end{multicols}


\section*{Άσκηση 15  \hfill \small{}}
Να υπολογίσετε την τιμή των παραστάσεων:
%\begin{multicols}{2}
\begin{enumerate}[1)]
 \item $Α=(α-1)(α+1)(α+2)(2α-5)$, όταν $α=3$
 \item $Β=(1-β)(2-β)(β+1)(2β-1)$, όταν $β=-1$
 \item $Γ=(3γ+7)(-3γ+1)-(9γ-γ)$, όταν $γ=\dfrac{1}{3}$
 \item $Δ=δ^{2}+δ^{3}+δ\cdot (1+2δ)$, όταν $δ=2$
 \item $Ε=ε(2ε-10)-(5-3ε)ε^{2}$, όταν $ε=10$
\end{enumerate}
%\end{multicols}


\section*{Άσκηση 16  \hfill \small{}}
Να λύσετε τις παρακάτω εξισώσεις:
%\begin{multicols}{2}
\begin{enumerate}[1)]
 \item $3+x=12$
 \item $7+2x=7$
 \item $x-5=2\cdot3$
 \item $10+10x=-100$
 \item $2\cdot x+3\cdot x-4\cdot x=8+2\cdot 4-4^{2}+5^{2}$
 \item $\dfrac{x}{2}=6+8-10$
 \item $\dfrac{3}{x}=7+\dfrac{1}{2}\cdot \dfrac{4}{3}$
 \item $3\cdot x-5\cdot x+2\cdot x=7\cdot (3+6+10+1)$
 \item $6x-10x=-30+6$
 \item $5x-10x=35$
\end{enumerate}
%\end{multicols}

%\section*{Άσκηση 16  \hfill \small{}}
%\begin{multicols}{2}
%\begin{enumerate}[1)]
% \item 
%\end{enumerate}
%\end{multicols}

\end{document}
